% =========================================================
% CAPITOLO 4 -- IMPLEMENTAZIONE
% =========================================================

\chapter{Implementazione}
\label{chap:implementazione}

Il Capitolo~\ref{chap:metodologia} ha formalizzato il \textit{framework} ibrido come un sistema di
costrutti matematici distinti per responsabilità e temporalità: un ipergrafo statico per la
semantica certificativa, un \gls{ATG} dinamico per la modellazione delle metriche a
\textit{runtime}, una funzione di \textit{mapping} bidirezionale e una \textit{pipeline} a quattro
fasi che trasforma la telemetria grezza in giudizi di \textit{compliance} predittivi. Il presente
capitolo descrive come questa architettura concettuale è tradotta in un sistema \textit{software}
operativo, illustrando le scelte ingegneristiche che governano l'organizzazione del codice, la
separazione delle responsabilità, la gestione della configurazione e l'implementazione dei singoli
moduli. L'attenzione è posta non solo su cosa il sistema fa, ma su come e perché determinate scelte
di \textit{design} siano state adottate in luogo di alternative disponibili.

% =========================================================
\section{Stack Tecnologico e Metodologia di Sviluppo}
\label{sec:stack_metodologia}
% =========================================================

Il sistema è interamente scritto in Python 3.11. La scelta di Python è motivata da due ragioni
convergenti. La prima è la disponibilità di un ecosistema maturo e consolidato per tutte le
aree funzionali richieste: analisi di serie temporali, \textit{graph processing}, apprendimento
automatico e test statistici. La seconda è la capacità del linguaggio di far coesistere nello
stesso progetto logiche algoritmiche di natura eterogenea, vale a dire elaborazione di grafi, analisi
statistica delle serie temporali, modelli predittivi e test causali, senza richiedere
ambienti di esecuzione separati o strati di integrazione tra componenti distinti.

Le dipendenze esterne si organizzano su quattro aree funzionali. Per la manipolazione vettorializzata
dei dati e delle serie temporali si impiegano pandas e NumPy: tutta la pipeline di trasformazione
delle metriche (calcoli differenziali su contatori cumulativi, gestione dei NaN, \textit{join}
sugli indici di \textit{timestamp}) è implementata tramite operazioni vettoriali che evitano
l'iterazione esplicita sulle righe, garantendo le prestazioni necessarie per elaborare finestre
temporali di diverse decine di migliaia di record. Per la modellazione topologica si utilizza
NetworkX, che fornisce un'implementazione nativa di grafi diretti (\textit{nx.DiGraph}) pienamente
compatibile con le operazioni di intersezione e interrogazione semantica richieste dal Layer~1.
La scelta di NetworkX rispetto ad alternative come graph-tool o igraph è motivata dalla sua
capacità di trattare i dizionari Python come attributi nativi di nodi e archi: questo rende il
campo \textit{hyperedges} di ogni arco un attributo di prima classe, interrogabile con la stessa
sintassi del grafo e modificabile senza operazioni di serializzazione, caratteristica essenziale
per l'implementazione della Strategia di Annotazione Semantica. Per il \textit{forecasting} si
integrano Prophet (Meta Research) per le serie con stagionalità, statsmodels con l'interfaccia
SARIMAX per i modelli \gls{ARIMA}, scikit-learn per la regressione lineare e per l'Isolation
Forest. Per l'analisi causale si utilizzano il modulo \textit{statsmodels.tsa.stattools} per i
test di Augmented Dickey-Fuller e di causalità di Granger, e scipy per il coefficiente di
Pearson.

Tutte le dipendenze sono dichiarate con versione fissata nel file \textit{pyproject.toml} e
riproducibili tramite \textit{uv}, garantendo che l'ambiente di esecuzione sia identico tra
sviluppo, integrazione continua e produzione.

Lo sviluppo ha seguito una metodologia \gls{TDD} in modo rigoroso: ogni modulo è stato costruito
a partire da una suite di test unitari che formalizzano il contratto dell'interfaccia pubblica
prima ancora che l'implementazione esista. Questa sequenza ha avuto conseguenze architetturali
dirette: alcune interfacce sono state ridisegnate più volte durante lo sviluppo per rendere i
test più semplici da costruire in isolamento, e ogni ridisegno ha prodotto un'interfaccia più
minimale e più facilmente componibile.

\subsection{Architettura della Repository}
\label{sec:architettura_repo}

La \textit{repository} adotta una struttura a pacchetti nidificati sotto la radice \textit{src/},
che riflette direttamente la stratificazione concettuale del \textit{framework}: le primitive
trasversali al livello più basso, i layer del modello al livello intermedio, le fasi della
\textit{pipeline} al livello superiore.

\begin{table}[htbp]
\centering
\caption{Organizzazione modulare del pacchetto \textit{src} con corrispondenza ai layer
metodologici.}
\label{tab:repo_structure}
\begin{tabular}{llll}
\toprule
\textbf{Pacchetto} & \textbf{Modulo} & \textbf{Layer / Fase} & \textbf{Responsabilità} \\
\midrule
\textit{src/utils/}   & \textit{logging\_setup.py}      & Infrastruttura & Primitiva di \textit{logging} uniforme \\
                      & \textit{config\_loader.py}      & Infrastruttura & Lettura e validazione YAML \\
\textit{src/ingestion/}     & \textit{converter.py}           & ETL            & Conversione dataset nel formato canonico \\
\textit{src/layer1/}  & \textit{topology\_builder.py}   & Layer~1        & Costruzione di $H_{\text{cert}}$ \\
\textit{src/layer2/}  & \textit{atg\_builder.py}        & Layer~2        & Sequenza di snapshot $G_t$ \\
                      & \textit{pbo\_builder.py}        & Layer~2        & PBO, PAS, norma di Frobenius \\
\textit{src/layer3/}  & \textit{feature\_selector.py}   & Layer~3 / Fase~I & Selezione di $M_{\Phi_i}$ dalla topologia \\
\textit{src/phase1/}  & \textit{stat\_forecaster.py}    & Fase~I           & Previsione locale delle metriche \\
\textit{src/phase2/}  & \textit{causal\_analyzer.py}    & Fase~II        & Analisi causale guidata dalla topologia \\
\textit{src/phase3/}  & \textit{structural\_monitor.py} & Fase~III       & Monitoraggio gerarchico a quattro livelli \\
\textit{src/phase4/}  & \textit{alert\_generator.py}    & Fase~IV        & Sintesi e classificazione dell'\textit{alert} \\
\bottomrule
\end{tabular}
\end{table}

Le dipendenze tra pacchetti scorrono esclusivamente verso il basso nella gerarchia: i moduli di
una fase non importano mai da una fase superiore, e i layer non dipendono dalle fasi. Ogni
cartella è responsabile di un solo costrutto del formalismo, e tutte le eccezioni, i messaggi di
\textit{log} e i parametri di configurazione di un modulo fanno riferimento esclusivamente agli
oggetti di quel costrutto. Questa organizzazione consente di verificare la logica del
TopologyBuilder in completo isolamento, senza istanziare alcun componente di previsione o di
\textit{anomaly detection}.

\subsection{Sistema di Configurazione}
\label{sec:config_system}

Il sistema di configurazione governa due aspetti fondamentali del \textit{framework}: la
separazione tra la struttura statica del sistema da monitorare e i parametri algoritmici dinamici
della \textit{pipeline}, e il meccanismo attraverso cui i moduli accedono a queste informazioni
in modo sicuro, coerente e verificabile. Le scelte adottate in questa area hanno implicazioni
dirette sulla testabilità di tutti i moduli e sulla manutenibilità a lungo termine del sistema.

\subsubsection{Struttura Topologica e Parametri}
\label{subsec:sep_yaml}

Un requisito architetturale fondamentale del \textit{framework} è la separazione completa tra
la struttura del sistema da monitorare e i parametri algoritmici della \textit{pipeline}. Queste
due categorie hanno cicli di vita radicalmente diversi: la topologia del sistema distribuito, le
definizioni dei \textit{compliance set} e gli \gls{SLA} variano su scala di mesi o di cicli di
certificazione; i parametri algoritmici, come le soglie Pearson, il fattore di smorzamento
\gls{EWMA} o i criteri di classificazione dell'\textit{alert}, variano a ogni sessione di
ottimizzazione. Mescolarli in una struttura monolitica renderebbe ogni sessione di \textit{tuning}
un rischio di corruzione accidentale della topologia.

Il sistema adotta pertanto due file di configurazione distinti. Il file \textit{config/topology.yaml}
formalizza la struttura statica del sistema da monitorare. Contiene le definizioni dei nodi e
degli archi del grafo di dipendenza tra i servizi, la specifica dei \textit{compliance set} con
il relativo tipo di topologia, le soglie \gls{SLA} e, per i set lineari, la sequenza del
percorso critico certificato. Vi sono inoltre le liste dei nomi canonici delle metriche di nodo e
di arco, i percorsi dei CSV canonici e i metadati del dataset, tra cui il parametro
\textit{window\_duration\_seconds} utilizzato come \textit{fallback} nel calcolo del
\textit{throughput}.

Il file \textit{config/pipeline\_params.yaml} raccoglie esclusivamente parametri algoritmici,
organizzati in cinque sezioni funzionali: i parametri del Probabilistic Behavioral Overlay,
comprendenti la metrica dei pesi della matrice stocastica e l'etichetta del \textit{gold
standard}; le impostazioni del \textit{forecasting}, con l'orizzonte temporale, il \textit{routing}
dei modelli, la configurazione ARIMA e la soglia di divergenza dalla \textit{baseline}; le soglie
per l'analisi causale (Pearson, Granger, Transfer Entropy); i parametri dell'\textit{anomaly
detection}, che includono z-score, Isolation Forest, \gls{CUSUM} e validatore strutturale; e
infine i criteri di classificazione degli \textit{alert}, con le soglie di \textit{lead time}
per i tre livelli di criticità.

\begin{table}[htbp]
\centering
\caption{Criterio di assegnazione delle informazioni ai due file di configurazione e motivazione
della separazione.}
\label{tab:config_sep}
\resizebox{\textwidth}{!}{%
\begin{tabular}{lll}
\toprule
\textbf{Categoria di informazione} & \textbf{File} & \textbf{Cambia quando} \\
\midrule
Identità dei nodi e ordine                 & \textit{topology.yaml}         & Cambia l'architettura \\
Definizione degli archi                     & \textit{topology.yaml}         & Cambia l'architettura \\
Membership dei \textit{compliance set}      & \textit{topology.yaml}         & Cambiano le proprietà certificate \\
Soglie \gls{SLA} per metrica                & \textit{topology.yaml}         & Cambiano i requisiti certificativi \\
Percorso critico e \textit{topology\_type}  & \textit{topology.yaml}         & Cambia l'architettura \\
\midrule
Soglie Pearson, Granger, Transfer Entropy   & \textit{pipeline\_params.yaml} & Si ricalibra l'analisi causale \\
Iperparametri Isolation Forest              & \textit{pipeline\_params.yaml} & Si ottimizza l'\textit{anomaly detection} \\
Fattori \gls{EWMA} e soglie \gls{CUSUM}    & \textit{pipeline\_params.yaml} & Si ricalibra il monitoraggio \\
Soglie \textit{lead time} (livelli colore)  & \textit{pipeline\_params.yaml} & Si aggiorna la \textit{policy} di allerta \\
Orizzonte di previsione $\tau$              & \textit{pipeline\_params.yaml} & Si adatta l'orizzonte operativo \\
Metriche non-lineari per \gls{LSTM}         & \textit{pipeline\_params.yaml} & Si aggiorna il \textit{routing} \\
\bottomrule
\end{tabular}%
}
\end{table}

Questa separazione ha implicazioni operative concrete. Il file \textit{topology.yaml} è
versionato insieme alla documentazione di certificazione del sistema distribuito: ogni modifica
alla membership di un \textit{compliance set} o a una soglia \gls{SLA} costituisce una modifica
al modello certificativo e dovrebbe essere accompagnata da una nuova sessione di validazione. Il file
\textit{pipeline\_params.yaml} può essere aggiornato a ogni sessione di ottimizzazione senza
coinvolgere il processo di certificazione, riducendo drasticamente il rischio che una sessione
di \textit{tuning} degli algoritmi introduca accidentalmente una variazione nella definizione
dei \textit{compliance set}. Un ulteriore beneficio è la testabilità per livelli: i test del
TopologyBuilder usano una versione di \textit{topology.yaml} ridotta con una topologia semplificata,
mentre i test delle fasi superiori usano un \textit{pipeline\_params.yaml} con soglie calibrate
per la copertura dei casi limite, senza mai alterare i parametri di produzione.

\subsubsection{ConfigLoader: Lettura Lazy e Validazione Eager}
\label{subsec:config_loader}

Il modulo ConfigLoader astrae l'accesso ai due file YAML con tre garanzie architetturali
distinte, che insieme garantiscono che il sistema non possa avviarsi in uno stato incompleto,
che le letture successive al disco siano evitate e che la mutazione esterna del dizionario
restituito non possa corrompere lo stato interno.

La prima garanzia è il caricamento \textit{lazy}: il file viene letto dal disco solo
alla prima chiamata del metodo corrispondente (\textit{load\_topology()} o
\textit{load\_pipeline\_params()}), e il risultato è mantenuto in \textit{cache} per le chiamate
successive. Questa scelta consente di istanziare ConfigLoader nel costruttore di qualsiasi modulo
senza accedere al \textit{filesystem} fino a quando i dati non sono effettivamente necessari,
migliorando la testabilità tramite \textit{mock} parziali.

La seconda garanzia è la validazione \textit{eager}: all'atto del primo accesso, il
metodo privato \textit{\_load\_and\_validate}, mostrato nel Codice~\ref{lst:load_validate},
verifica la presenza di tutte le chiavi obbligatorie
al livello radice, interrompendo con \textit{ValueError} alla prima chiave mancante e includendo
il nome esatto della chiave nel messaggio d'errore. Questo anticipa i fallimenti alla fase di
inizializzazione del sistema, non all'esecuzione, evitando la classe di errori in cui un componente
si avvia silenziosamente con una configurazione incompleta e fallisce molto più tardi con un errore
opaco.

\begin{lstlisting}[
  language=Python,
  caption={Metodo \_load\_and\_validate in src/utils/config\_loader.py.},
  label={lst:load_validate}
]
@staticmethod
def _load_and_validate(
    path: Path, required_keys: tuple[str, ...]
) -> dict:
    if not path.exists():
        raise FileNotFoundError(
            f"File di configurazione non trovato: {path.resolve()}"
        )
    try:
        with path.open(encoding="utf-8") as fh:
            data: dict = yaml.safe_load(fh)
    except yaml.YAMLError as exc:
        raise ValueError(
            f"{path.name} contiene YAML sintaticamente "
            f"non valido: {exc}"
        ) from exc
    if not isinstance(data, dict):
        raise ValueError(
            f"{path.name} non contiene un mapping YAML valido "
            f"(ottenuto: {type(data).__name__})"
        )
    for key in required_keys:
        if key not in data:
            raise ValueError(
                f"Chiave obbligatoria mancante in "
                f"{path.name}: '{key}'"
            )
    return data
\end{lstlisting}

La terza garanzia è l'isolamento della \textit{cache}: ogni chiamata restituisce una
copia profonda del dizionario tramite \textit{copy.deepcopy}. Senza questo meccanismo, qualsiasi
modulo che ricevesse il dizionario e lo modificasse corromperebbe la struttura condivisa,
producendo comportamenti non deterministici nelle chiamate successive. La suite di test verifica
questo isolamento esplicitamente anche per strutture annidate, come mostra il Codice~\ref{lst:cache_isolation}:

\begin{lstlisting}[
  language=Python,
  caption={Test di isolamento della cache in tests/test\_config\_loader.py.},
  label={lst:cache_isolation}
]
def test_load_topology_cache_deep_isolation(tmp_path):
    mock_loader = ConfigLoader(topo, pipe)
    t1 = mock_loader.load_topology()
    original_len = len(
        t1["compliance_sets"]["CS_test"]["nodes"]
    )
    t1["compliance_sets"]["CS_test"]["nodes"].append(
        "injected-sentinel"
    )
    t2 = mock_loader.load_topology()
    assert len(
        t2["compliance_sets"]["CS_test"]["nodes"]
    ) == original_len
\end{lstlisting}

Il file \textit{topology.yaml} richiede una validazione aggiuntiva rispetto alla semplice
presenza delle chiavi radice: il campo \textit{metadata.window\_duration\_seconds} è obbligatorio
e deve essere un numero positivo, perché viene usato dal convertitore come \textit{fallback}
per il calcolo del \textit{throughput} nelle finestre con $\Delta t$ non calcolabile. La
sua assenza o un valore non positivo produce un \textit{ValueError} esplicito al momento del
caricamento, prevenendo fallimenti silenziosi nella fase di conversione.

\subsubsection{La Primitiva di Logging}
\label{subsec:logging_setup}

La classe LoggingSetup fornisce una primitiva trasversale che garantisce un formato di
\textit{logging} uniforme in tutti i moduli del \textit{framework}. Ogni modulo ottiene un
\textit{logger} nominato con il proprio percorso Python completo, rendendo possibile il
filtraggio per componente durante l'esecuzione in produzione e durante le sessioni di
\textit{debugging}. La primitiva è idempotente per costruzione: se un \textit{logger} con
lo stesso nome è già stato configurato, la chiamata successiva restituisce lo stesso
\textit{logger} senza aggiungere \textit{handler} duplicati. Senza questa garanzia, in ambienti
dove i moduli vengono importati più volte, ad esempio nei test con ricaricamento dinamico,
ogni messaggio verrebbe emesso più volte, producendo \textit{output} difficile da interpretare.

% =========================================================
\section{Il Formato Canonico di Input}
\label{sec:etl_converter}
% =========================================================

La \textit{pipeline} del \textit{framework} non opera mai direttamente sui dati grezzi
prodotti dagli strumenti di osservabilità. Tra i dati grezzi e i moduli di analisi si
interpone uno strato di conversione il cui scopo è produrre un \textbf{formato canonico
intermedio} indipendente dal dataset specifico. Questa separazione consente di sostituire
la sorgente dati senza modificare una riga di codice al di fuori del pacchetto
\textit{src/ingestion/}.

Il formato canonico si compone di tre file CSV con schemi fissi, che rispecchiano direttamente
la distinzione strutturale del Layer~2 tra \textit{feature} persistenti di nodo, \textit{feature}
effimere di arco e informazione di \textit{ground truth}. Il file \textit{node\_metrics.csv} è
organizzato in formato \textit{long}: ogni riga rappresenta il vettore di stato di un singolo
microservizio in una singola finestra temporale, con \textit{timestamp} in microsecondi, il
nome del servizio e le metriche di stato interno dichiarate in
\textit{topology["node\_metrics"]}. La chiave composita di unicità è la coppia
\textit{(timestamp, node\_id)}. Il file \textit{edge\_metrics.csv} veicola le
\textit{feature} effimere degli archi, con colonne \textit{timestamp}, \textit{edge\_id},
sorgente, destinazione e le metriche di qualità dichiarate in
\textit{topology["edge\_metrics"]}, con chiave \textit{(timestamp, edge\_id)}. Il file
\textit{ground\_truth.csv} associa a ogni finestra temporale la \textit{label} binaria di
anomalia, la lista JSON dei nodi anomali e i metadati dell'esperimento, con
\textit{timestamp} come chiave.

Qualunque componente di conversione che produca questi tre file con gli schemi descritti è
pertanto compatibile con l'intera \textit{pipeline} del \textit{framework},
e rappresenta l'unico modulo \textit{dataset-specific}
dell'intera architettura: si occupa di estrarre i metadati, aggregare le tracce per finestra
temporale, calcolare le metriche con le opportune correzioni sui contatori cumulativi e
produrre il \textit{ground truth}. L'implementazione specifica adottata nella valutazione
sperimentale è descritta nel Capitolo~\ref{chap:sperimentazione}.

% =========================================================
\section{Layer 1: L'Ipergrafo di Certificazione}
\label{sec:impl_layer1}
% =========================================================

Il Layer~1 del \textit{framework} è responsabile della rappresentazione statica della struttura
certificativa del sistema: quali componenti appartengono a quale proprietà certificata, quali
risorse sono condivise tra proprietà diverse, e quale sequenza costituisce il percorso critico
certificato per ogni proprietà lineare. Questa informazione è codificata nel file
\textit{topology.yaml} e tradotta in strutture dati interrogabili dal modulo TopologyBuilder,
che fornisce a tutti i moduli superiori un'interfaccia uniforme e tipizzata per accedere alla
topologia del sistema.

\subsection{Annotazione Semantica e Scelta di NetworkX}
\label{subsec:annotazione_semantica}

Il TopologyBuilder traduce il file \textit{topology.yaml} nella struttura dati che implementa
il formalismo $H_{\text{cert}} = (V, \{H_{\Phi_1}, \ldots, H_{\Phi_m}\})$. La sfida
ingegneristica principale è che NetworkX non dispone di un tipo di dato ipergrafale nativo:
non esiste una classe \textit{Hypergraph} che modelli direttamente la semantica di un iperarco
come insieme di nodi con logica AND atomica. Si presentano pertanto due strategie di
implementazione, analizzate anche nel Capitolo~\ref{chap:metodologia}: l'Espansione Topologica
e l'Annotazione Semantica.

L'Espansione Topologica introduce nodi virtuali di sincronizzazione, uno per ogni
\textit{compliance set}, collegati tramite archi fittizi a tutti i nodi del set. Questa strategia
rende la semantica AND esplicita nella topologia, ma introduce $m$ nodi aggiuntivi (uno per
proprietà) le cui \textit{feature} devono essere ricalcolate a ogni passo temporale, producendo
un \textit{overhead} computazionale crescente con il numero di proprietà monitorate. Tale strategia produce
inoltre strutture eterogenee, con nodi reali e nodi fittizi nello stesso grafo, che complicano le
operazioni di interrogazione e la compatibilità con le \textit{pipeline} di osservabilità
industriali, le quali producono metriche associate agli archi e ai nodi reali del sistema.

La Strategia di Annotazione Semantica, adottata nell'implementazione, mantiene invariata la
topologia reale del sistema da monitorare e codifica la semantica dei \textit{compliance set}
come attributo \textit{hyperedges} su ogni arco del grafo: una lista dei nomi dei
\textit{compliance set} a cui l'arco appartiene secondo la definizione formale
$A(H_{\Phi_i}) = \{e = (u,v) \mid u \in H_{\Phi_i} \wedge v \in H_{\Phi_i}\}$. Questo
approccio è compatibile con qualsiasi libreria di analisi su grafi che supporti attributi sugli
archi, non introduce nodi fittizi, non richiede ricalcolo a \textit{runtime} e consente di
aggiungere una nuova proprietà certificata aggiornando solo le annotazioni degli archi già
esistenti. La semantica AND del \textit{compliance set} è rispettata integralmente dalla
logica applicativa dei moduli superiori (FeatureSelector, StructuralMonitor, CausalAnalyzer)
che interrogano questo attributo: l'assenza di un tipo \textit{Hypergraph} nativo non implica
una simulazione approssimata, ma una codifica equivalente della stessa semantica in una
struttura dati differente.

La scelta specifica di NetworkX è motivata dal suo modello a dizionari Python nativi per
gli attributi: \textit{g.add\_edge(u, v, hyperedges=[...])} rende il campo \textit{hyperedges}
interrogabile con la stessa sintassi del grafo e modificabile senza operazioni di
serializzazione o cast di tipo. Questa flessibilità, assente in librerie come igraph (che usa
schemi fissi per gli attributi), è fondamentale per mantenere l'annotazione semantica
aggiornabile in risposta a cambiamenti della configurazione.

\subsection{Il Metodo build() e il Grafo Annotato}
\label{subsec:build_method}

La costruzione del grafo avviene tramite il metodo \textit{build()}, mostrato nel
Codice~\ref{lst:build}, che implementa un pattern di \textit{caching} idempotente:
alla prima chiamata il grafo viene costruito e memorizzato
internamente; alle chiamate successive viene restituita una copia profonda del grafo già
costruito, senza ricalcolare. L'algoritmo di annotazione sfrutta la precalcolazione nel
costruttore: i set di nodi di ogni \textit{compliance set} vengono convertiti in strutture
\textit{set[str]} e memorizzati in \textit{\_cs\_node\_sets}, in modo che il \textit{lookup}
dell'appartenenza di un nodo a un \textit{compliance set} abbia costo $O(1)$ anziché $O(|V|)$
per ogni arco.

\begin{lstlisting}[
  language=Python,
  caption={Metodo build() in src/layer1/topology\_builder.py.},
  label={lst:build}
]
def build(self) -> nx.DiGraph:
    if self._graph is not None:
        return copy.deepcopy(self._graph)

    g = nx.DiGraph()
    for node in self._topology["nodes"]:
        g.add_node(node["id"])

    declared_nodes = {n["id"] for n in self._topology["nodes"]}
    for edge in self._topology["edges"]:
        for endpoint_key in ("source", "target"):
            ep = edge[endpoint_key]
            if ep not in declared_nodes:
                raise ValueError(
                    f"Arco '{edge['id']}': endpoint '{ep}' "
                    "non presente in topology['nodes']. "
                    "Verifica topology.yaml."
                )

    for edge in self._topology["edges"]:
        u, v = edge["source"], edge["target"]
        hyperedges = [
            name
            for name, node_set in self._cs_node_sets.items()
            if u in node_set and v in node_set
        ]
        g.add_edge(u, v, id=edge["id"], hyperedges=hyperedges)

    covered: set[str] = set()
    for cs in self._topology["compliance_sets"].values():
        covered.update(cs.get("nodes", []))
    isolated = {n["id"] for n in self._topology["nodes"]} - covered
    if isolated:
        logger.warning(
            "Nodi in V non appartenenti ad alcun compliance set "
            "(blind spot di monitoraggio): %s",
            sorted(isolated),
        )

    self._graph = g
    return copy.deepcopy(self._graph)
\end{lstlisting}

La validazione degli \textit{endpoint} è separata dalla costruzione degli archi: il primo
ciclo controlla tutti gli archi prima che qualsiasi arco sia aggiunto al grafo, in modo da
garantire che un errore di configurazione produca un'eccezione con il nome dell'arco
problematico, senza lasciare il grafo in uno stato parzialmente costruito.

Il ciclo di annotazione semantica include una seconda verifica, attivata dopo l'assegnazione
di \textit{hyperedges}: se l'attributo risulta una lista vuota e tuttavia entrambi gli
\textit{endpoint} appartengono ad almeno un \textit{compliance set}, il metodo emette
\textit{logger.warning}. Questo scenario si verifica quando i due estremi di un arco
appartengono a \textit{compliance set} distinti e non condivisi: l'arco non rientra in
$A(H_{\Phi_i})$ per nessuna proprietà e non è raggiungibile nemmeno da
$M^{\text{interf}}$, risultando strutturalmente invisibile all'intero \textit{framework}.
Il warning segnala una potenziale lacuna di copertura del monitoraggio prima che qualsiasi
analisi sia eseguita, evitando che la situazione rimanga silente fino all'esecuzione.

La chiamata finale \textit{copy.deepcopy(self.\_graph)} implementa il \textit{Pattern Prototype}: il
grafo originale è conservato come \textit{template} immutabile internamente; ogni consumatore
riceve una copia indipendente e può modificarla liberamente senza corrompere la struttura
centralizzata. Questa garanzia è verificata dal Codice~\ref{lst:test_build_copy}:

\begin{lstlisting}[
  language=Python,
  caption={Test test\_build\_returns\_independent\_copy in tests/test\_topology\_builder.py.},
  label={lst:test_build_copy}
]
def test_build_returns_independent_copy(builder):
    g1 = builder.build()
    g1.add_node("injected-sentinel")
    g2 = builder.build()
    assert "injected-sentinel" not in g2.nodes()
\end{lstlisting}

\subsection{L'Interfaccia pubblica e le Query Semantiche}
\label{subsec:topology_queries}

L'interfaccia pubblica del TopologyBuilder espone sette metodi che traducono direttamente le
operazioni formali del Capitolo~\ref{chap:metodologia} in API Python tipizzate. Tutti i moduli
superiori (FeatureSelector, StructuralMonitor, CausalAnalyzer) accedono alla topologia
esclusivamente attraverso questi metodi, senza mai interrogare il grafo NetworkX direttamente.
Questa scelta segue il principio di Law of Demeter e consente di sostituire internamente la
rappresentazione del grafo senza modificare alcun modulo dipendente.

Il metodo \textit{get\_edges\_for\_compliance\_set} implementa $A(H_{\Phi_i})$, calcolando
dinamicamente la condizione $u \in H_{\Phi_i} \wedge v \in H_{\Phi_i}$ su ogni arco del grafo.
Il risultato non è memorizzato in \textit{cache} perché la chiamata è $O(|E|)$ e il numero
di archi in un sistema distribuito tipico è sufficientemente ridotto da rendere il ricalcolo
trascurabile.

Il metodo \textit{get\_interference\_edges} implementa $M^{\text{interf}}$, restituendo gli
archi $e = (u, v)$ tali che $v \in \text{Shared}(H_{\Phi_i}, H_{\Phi_j})$ e
$u \notin H_{\Phi_i}$. La funzione solleva \textit{ValueError} se \textit{target\_cs} e
\textit{other\_cs} coincidono, perché l'auto-interferenza non è definita semanticamente. Il
risultato può essere l'insieme vuoto: questo non è un caso degenere ma una proprietà strutturale
legittima di topologie in cui tutti gli archi entranti verso i nodi condivisi provengono da nodi
già interni al \textit{compliance set} di destinazione.

Il metodo \textit{get\_critical\_path} include una doppia validazione strutturale: verifica
che ogni nodo della sequenza appartenga al \textit{compliance set} e che ogni coppia
consecutiva corrisponda a un arco reale nella topologia, sollevando \textit{ValueError} in
entrambi i casi. Questa validazione \textit{eager}, eseguita alla prima chiamata, anticipa i
fallimenti alla fase di configurazione e non all'esecuzione, dove un percorso mal formato
produrrebbe un Path Adherence Score semanticamente privo di significato. Per i
\textit{compliance set} di tipo \textit{parallel}, il metodo restituisce una lista vuota,
evitando la necessità di una condizione speciale nei chiamanti. Qualora il
\textit{topology\_type} non sia né \textit{linear} né \textit{parallel}, il metodo emette un
\textit{warning} con il valore non riconosciuto e restituisce una lista vuota, senza sollevare
eccezioni: questa scelta di \textit{fail-soft} consente al sistema di continuare a operare
anche in presenza di configurazioni parzialmente valide.

La suite di test del TopologyBuilder comprende 38 casi che verificano le proprietà
dell'interfaccia pubblica: struttura del grafo, annotazione semantica, nodi condivisi,
$A(H_{\Phi_i})$, $M^{\text{interf}}$, idempotenza di \textit{build()} e isolamento della
\textit{cache}, nonché tutti i comportamenti d'errore attesi per configurazioni malformate.

% =========================================================
\section{Layer 2: ATG e Probabilistic Behavioral Overlay}
\label{sec:impl_layer2}
% =========================================================

Il Layer~2 si occupa di due compiti concettualmente distinti ma strettamente complementari. Il
primo è costruire la sequenza temporale di snapshot $G_t$ a partire dai tre CSV canonici,
allineando e deduplicando i dati provenienti da sorgenti con \textit{timestamp} potenzialmente
non coincidenti. Il secondo è trasformare quella sequenza in una rappresentazione comportamentale
stocastica che cattura non i valori puntuali delle metriche, ma la distribuzione del traffico tra
i percorsi del sistema. Questi due compiti sono assegnati a due classi distinte con responsabilità
non sovrapposte: ATGBuilder e PBOBuilder. La separazione è una scelta di \textit{design}
deliberata motivata da ragioni di testabilità e di coesione: i 32 test dell'ATGBuilder non
istanziano mai un PBOBuilder, e i 27 test del PBOBuilder non leggono mai un CSV reale.

\subsection{ATGBuilder: Snapshot e Allineamento Temporale}
\label{subsec:impl_atg}

L'ATGBuilder implementa il Layer~2 del \textit{framework} costruendo la sequenza temporale di
snapshot $G_t = (V, E_t, \mathbf{X}_{V,t}, \mathbf{X}_{E,t})$ a partire dai tre CSV canonici.
Ogni snapshot è rappresentato come un dizionario Python con sei campi che corrispondono
esattamente agli elementi della quadrupla formale. Il campo \textit{timestamp} è un intero in
microsecondi, non convertito in \textit{datetime64}: la conversione è delegata ai soli moduli
di \textit{forecasting} che la richiedono (Prophet), evitando una trasformazione globale non
necessaria per gli altri modelli. Il campo \textit{nodes} è una mappa da \textit{node\_id} al
dizionario delle metriche di nodo; il campo \textit{edges} una mappa da \textit{edge\_id} al
dizionario delle metriche di arco più i campi \textit{source} e \textit{target}. I campi
\textit{label}, \textit{anomaly\_type} e \textit{anomaly\_node\_ids} rispecchiano il
\textit{ground truth} della finestra.

Il metodo \textit{build()} accetta DataFrame opzionali in memoria (usati nei test) e, se assenti,
li legge dai percorsi CSV passati al costruttore. Questo schema a iniezione di dipendenza consente
ai test di fornire dati sintetici precisi senza richiedere alcun file su disco. Ogni chiamata a
\textit{build()} esegue l'intera pipeline di costruzione: non è idempotente rispetto alla lettura
da disco, perché i dati sul disco potrebbero cambiare tra chiamate.

La pipeline interna si articola in tre fasi sequenziali. La prima fase è la deduplicazione dei
tre CSV. Prima di qualsiasi allineamento, il metodo applica tre deduplicazioni con chiavi distinte:
per \textit{timestamp} sul \textit{ground truth}, per coppia \textit{(timestamp, node\_id)} sui
nodi e per coppia \textit{(timestamp, edge\_id)} sugli archi. La strategia \textit{keep="first"}
è deliberata: si mantiene la prima occorrenza e si scarta la seconda senza tentare un \textit{merge}.
Un \textit{merge} produrrebbe valori metrici sintetici non presenti nel dataset originale,
violando il principio che ogni valore nell'output deve essere direttamente tracciabile a un record
del dataset grezzo. Le deduplicazioni possono eliminare una percentuale non trascurabile di
osservazioni: per questo il comportamento è documentato come limitazione strutturale del dataset
di addestramento, non come difetto del modulo.

La seconda fase è l'allineamento per \textit{inner join}. Il metodo calcola l'intersezione dei
\textit{timestamp} unici dei tre DataFrame deduplicati e include nel risultato finale solo le
finestre presenti simultaneamente in tutti e tre. Il metodo pre-indicizza i DataFrame per \textit{timestamp}
tramite \textit{groupby} prima del ciclo principale, riducendo il costo di \textit{lookup} da $O(n)$
a $O(1)$ per snapshot. Se l'intersezione è vuota, viene sollevato \textit{ValueError}: questa è
l'unica condizione che blocca la costruzione, perché un insieme vuoto di snapshot è
strutturalmente inutilizzabile per qualsiasi analisi successiva. Una piccola differenza nel numero
di \textit{timestamp} tra \textit{node\_metrics} e \textit{edge\_metrics} è il comportamento
atteso quando il convertitore scarta la prima finestra per nodo (per impossibilità di calcolare
il delta del contatore) ma non per gli archi: il join \textit{inner} la esclude automaticamente.

La terza fase è la verifica consultiva della consistenza tra i nodi e gli archi presenti nei CSV
e quelli dichiarati in \textit{topology.yaml}. Le discrepanze non bloccano la costruzione: nodi o
archi extra emettono \textit{logger.warning} con il nome dell'identificatore sconosciuto, e nodi
o archi mancanti emettono anch'essi \textit{warning}. Una verifica analoga viene eseguita sulle
colonne di \textit{node\_metrics.csv}: se una metrica dichiarata in
\textit{topology["node\_metrics"]} è strutturalmente assente dal file CSV (non solo
sporadicamente NaN, ma completamente mancante come colonna), viene emesso un avviso esplicito che
elenca le dimensioni interessate. Gli snapshot verranno costruiti con NaN su quelle dimensioni, e
i moduli di \textit{forecasting} a valle li gestiranno tramite \textit{dropna()} prima del
\textit{fitting}. Questa politica di \textit{fail-soft} è motivata dal fatto che le discrepanze
possono essere transitorie, ad esempio per un nodo temporaneamente assente durante un
\textit{rolling restart}, e non devono interrompere il monitoraggio.

La gestione del campo \textit{anomaly\_type} richiede attenzione particolare. Il
\textit{fault\_type} estratto dal nome file vale per l'intera sessione sperimentale, incluse le
finestre nominali all'inizio dell'esperimento. Per preservare la semantica corretta, il campo
\textit{anomaly\_type} è impostato a \textit{None} quando \textit{label == 0},
indipendentemente dal valore di \textit{fault\_type}: una finestra nominale è nominale per
definizione, anche se appartiene a una sessione di guasto. I valori \textit{fault\_type} pari a
NaN (ad esempio da esperimenti con metadati incompleti) sono convertiti a \textit{None} e non
alla stringa \textit{"nan"}: un test dedicato (\textit{test\_fault\_type\_nan\_produces\_none\_anomaly\_type})
verifica questo comportamento. Il campo \textit{anomaly\_node\_ids} con JSON malformato produce
una lista vuota con \textit{logger.warning}, senza propagare eccezioni ai moduli superiori.

I metodi ausiliari \textit{get\_node\_feature\_matrix} e \textit{get\_edge\_feature\_matrix}
materializzano le matrici $\mathbf{X}_{V,t}$ e $\mathbf{X}_{E,t}$ come DataFrame pandas
indicizzati per \textit{timestamp}, con le metriche come colonne. Questi DataFrame sono
direttamente impiegabili dal FeatureSelector senza ulteriori trasformazioni. Qualora
l'identificatore richiesto non sia presente in nessuno snapshot della lista, entrambi i
metodi restituiscono un DataFrame vuoto ed emettono \textit{logger.warning} con il nome
dell'identificatore. Questo comportamento è coerente con quello di \textit{build()}, che registra attivamente gli
identificatori non riconosciuti durante la fase di costruzione, garantendo la visibilità
dell'assenza in entrambe le direzioni. I metodi statici
\textit{get\_nominal\_snapshots} e \textit{get\_anomalous\_snapshots} filtrano la lista per
etichetta e, il secondo, opzionalmente per tipo di anomalia con corrispondenza esatta sulla
stringa: \textit{"cpu"} non è una \textit{substring} di \textit{"cpu\_mem"}.

\subsection{PBOBuilder: Matrice Stocastica e Monitoraggio Comportamentale}
\label{subsec:impl_pbo}

Il PBOBuilder implementa il livello \gls{PBO} $G_{\text{Behavior}}(t) = (V, E_{\text{all}}, W_t)$,
trasformando gli snapshot dell'ATG in una rappresentazione che cattura non quanto il sistema sia
lento o saturato, ma come si distribuisce il traffico tra i percorsi disponibili. Questa
distinzione è il motivo per cui le due classi sono separate: l'ATGBuilder risponde alla domanda
\textit{quanto} (quanta latenza, quanta memoria, quanto carico), il PBOBuilder risponde alla
domanda \textit{come} (il traffico segue i percorsi certificati, o si sta redistribuendo verso
percorsi di \textit{fallback}?). Un sistema può apparire metricamente conforme mentre abbandona
progressivamente il percorso certificato: questa deviazione è visibile nel PBO prima ancora che
le soglie di latenza vengano superate.

Il costruttore del PBOBuilder valida subito che la metrica configurata in
\textit{pipeline\_params["pbo"]["weight\_metric"]} (di default \textit{throughput\_rps}) sia
presente nella lista \textit{edge\_metrics} di \textit{topology.yaml}. In caso contrario,
viene sollevato \textit{ValueError} al momento dell'istanziazione: questa scelta anticipa il
fallimento rispetto a una scoperta tardiva durante il calcolo dei pesi.

Il metodo \textit{compute\_transition\_weights} calcola la matrice stocastica $W_t$ per ogni
snapshot. Per ogni nodo con archi uscenti, normalizza il valore della metrica scelta su ciascun
arco rispetto alla somma totale degli archi uscenti dallo stesso nodo. Il meccanismo di
\textit{fallback} uniforme, mostrato nel Codice~\ref{lst:pbo_weights}, è uno degli aspetti
più critici dell'implementazione:

\begin{lstlisting}[
  language=Python,
  caption={Calcolo dei pesi di transizione in src/layer2/pbo\_builder.py.},
  label={lst:pbo_weights}
]
for src, edge_ids in self._source_to_edges.items():
    available = [eid for eid in edge_ids if eid in snap["edges"]]
    if not available:
        continue
    raw_values: dict[str, float] = {}
    for eid in available:
        try:
            raw_values[eid] = float(
                snap["edges"][eid][self._weight_metric]
            )
        except (TypeError, ValueError):
            raw_values[eid] = float("nan")

    has_invalid = any(
        not math.isfinite(v) or v < 0
        for v in raw_values.values()
    )
    total_tp = sum(
        v for v in raw_values.values()
        if math.isfinite(v) and v >= 0
    )
    if has_invalid or total_tp <= 0:
        w_uniform = 1.0 / len(available)
        for eid in available:
            weights[eid] = w_uniform
        logger.warning(
            "Throughput non valido per nodo '%s' al ts=%d: "
            "fallback uniforme 1/%d.",
            src, snap["timestamp"], len(available),
        )
    else:
        for eid in available:
            weights[eid] = raw_values[eid] / total_tp
\end{lstlisting}

Il ciclo esterno itera sui nodi con archi uscenti; per ciascuno, raccoglie i valori della metrica
di peso per gli archi uscenti presenti nello snapshot corrente. La condizione di \textit{fallback}
scatta se la somma è zero o non positiva, oppure se almeno un valore è NaN o negativo. In quel
caso, tutti i pesi del nodo vengono impostati a $1/n$ dove $n$ è il numero di archi uscenti,
preservando la proprietà stocastica. Il \textit{logger.warning} documenta il \textit{timestamp}
e il nome del nodo coinvolto, rendendo l'evento tracciabile in produzione.

Il metodo \textit{compute\_gold\_standard} calcola $W^{\text{gold}}$ come media dei pesi sulle
sole finestre con etichetta corrispondente al parametro \textit{gold\_standard\_label} (di default
0, ossia le finestre nominali). La \textit{label} di riferimento è letta dal YAML e non è mai
fissa nel codice. Il \textit{gold standard} copre tutti gli archi della topologia
(\textit{E\_all}): se un arco è assente in uno snapshot nominale (ad esempio perché il suo
\textit{throughput} è zero e il suo peso non è stato registrato), contribuisce con peso 0 alla
media. Un test dedicato verifica questa proprietà: costruire il \textit{gold standard} su due
snapshot nominali dove il primo non ha \textit{e6} produce
$W^{\text{gold}}[\text{e6}] = (0.0 + 0.6) / 2 = 0.3$, anziché il solo valore del secondo snapshot.

Il metodo \textit{compute\_path\_adherence} è applicabile solo ai \textit{compliance set} con
\textit{topology\_type == "linear"} e solleva \textit{ValueError} per i tipi paralleli. Come
mostrato nel Codice~\ref{lst:pas_frobenius}, il \gls{PAS} è implementato come prodotto cumulato
dei pesi degli archi lungo il percorso critico: se un arco del percorso certificato non è
presente nel dizionario dei pesi correnti, il suo contributo è zero, implementando la semantica
che un arco certificato senza traffico corrisponde a un percorso parzialmente abbandonato.

\begin{lstlisting}[
  language=Python,
  caption={Calcolo del PAS e della norma di Frobenius in src/layer2/pbo\_builder.py.},
  label={lst:pas_frobenius}
]
# Path Adherence Score
for entry in weight_series:
    weights = entry["weights"]
    pas = 1.0
    for eid in path_edge_ids:
        pas *= weights.get(eid, 0.0)
    result.append({"timestamp": entry["timestamp"], "pas": pas})

# Norma di Frobenius
all_eids = [e["id"] for e in self._edges]
for entry in weight_series:
    weights = entry["weights"]
    sq_sum = sum(
        (weights.get(eid, 0.0) - gold_standard.get(eid, 0.0)) ** 2
        for eid in all_eids
    )
    result.append({
        "timestamp": entry["timestamp"],
        "frobenius": math.sqrt(sq_sum),
    })
\end{lstlisting}

Il primo blocco implementa il \gls{PAS}: \textit{pas} inizializzato a 1.0 viene moltiplicato
per il peso di ogni arco del percorso critico. Il secondo blocco implementa la norma di
Frobenius: per ogni arco della topologia completa $E_{\text{all}}$, calcola il quadrato della
differenza tra il peso corrente e il peso del \textit{gold standard}, usando 0.0 come valore di
\textit{default} per entrambi i dizionari. Questo garantisce che un arco con peso 0.0 esplicito
e un arco assente dal dizionario producano identico contributo alla norma, proprietà verificata
da \textit{test\_frobenius\_explicit\_zero\_weight\_contributes}.

Il PBOBuilder è utilizzato dallo StructuralMonitor in due modalità distinte a seconda del tipo di
topologia del \textit{compliance set}: per le topologie lineari il monitoraggio opera sul \gls{PAS},
accumulando i suoi decrementi nel \gls{CUSUM}; per le topologie parallele, dove il \gls{PAS} non
è definito, il monitoraggio opera sulla norma di Frobenius, accumulando i suoi incrementi. La
scelta della metrica appropriata avviene a \textit{runtime} leggendo il campo
\textit{topology\_type} dalla configurazione.

% =========================================================
\section{Implementazione del Layer 3: Feature Selector}
\label{sec:impl_feature_selector}
% =========================================================

Il FeatureSelector implementa il Layer~3 del \textit{framework}, ossia la funzione di
\textit{mapping} $\mathcal{M}$ che traduce la specifica di un \textit{compliance set}
nell'insieme di serie temporali $M_{\Phi_i} = M^{\text{direct}} \cup M^{\text{interf}}$
pronte per il \textit{forecasting}. Il modulo non interroga mai direttamente il grafo NetworkX
e non accede mai ai CSV: delega tutte le operazioni topologiche al TopologyBuilder e riceve
i valori delle metriche dagli snapshot prodotti dall'ATGBuilder. Questa architettura segue il
principio \textit{Don't Repeat Yourself}: la logica di calcolo di $A(H_{\Phi_i})$ e di
$M^{\text{interf}}$ è implementata una sola volta nel TopologyBuilder; il FeatureSelector ne
consuma i risultati. Se il criterio di appartenenza di un arco a un \textit{compliance set}
cambiasse, la modifica dovrebbe essere applicata in un solo posto.

\subsection{Convenzione di Denominazione e Struttura di select\_features}
\label{subsec:feature_naming}

Il metodo principale \textit{select\_features} restituisce un dizionario che mappa stringhe
chiave a DataFrame pandas. La convenzione di denominazione delle chiavi è tripartita e
deterministica: le \textit{feature} di nodo seguono il formato
\textit{node:<node\_id>:<metrica>}, quelle di arco \textit{edge:<edge\_id>:<metrica>} e quelle
di interferenza \textit{interf:<edge\_id>:throughput\_rps}. Questa convenzione rende le chiavi
interpretabili sia per la logica del \textit{framework} (il CausalAnalyzer può ricostruire
l'appartenenza di una metrica a un nodo o a un arco analizzando il prefisso) sia per l'operatore
che legge i \textit{log}.

L'ordine di costruzione delle chiavi è determinato da regole precise. Le chiavi \textit{node:}
seguono \textit{sorted(cs\_nodes)}: l'ordine lessicografico sui nodi è fondamentale per la
riproducibilità del vettore di stato aggregato dello StructuralMonitor, che deve avere la stessa
struttura tra il \textit{fit()} e ogni successiva chiamata a \textit{monitor()}. Le chiavi
\textit{edge:} seguono l'ordine di dichiarazione degli archi in \textit{topology.yaml}. Le
chiavi \textit{interf:} seguono l'ordine di visita degli altri \textit{compliance set}.
Il Codice~\ref{lst:select_features} mostra la struttura completa del metodo.

\begin{lstlisting}[
  language=Python,
  caption={Metodo select\_features in src/layer3/feature\_selector.py.},
  label={lst:select_features}
]
def select_features(self, compliance_set_name, snapshots):
    cs_nodes = self._tb.get_compliance_set_nodes(compliance_set_name)
    internal_edges = self._tb.get_edges_for_compliance_set(
        compliance_set_name
    )
    interf_edges = self._collect_interference_edges(compliance_set_name)
    result = {}

    # M_direct -- feature di nodo (ordine lessicografico deterministico)
    for node_id in sorted(cs_nodes):
        for metric in self._node_metrics:
            key = f"node:{node_id}:{metric}"
            result[key] = self._build_node_series(
                node_id, metric, snapshots
            )

    # M_direct -- feature di arco (ordine topology.yaml)
    for src, tgt in internal_edges:
        edge_id = self._edge_id_lookup.get((src, tgt))
        if edge_id is None:
            continue
        for metric in self._edge_metrics:
            key = f"edge:{edge_id}:{metric}"
            result[key] = self._build_edge_series(
                edge_id, metric, snapshots
            )

    # M_interf -- solo throughput_rps, con deduplicazione tramite seen-set
    for src, tgt in interf_edges:
        edge_id = self._edge_id_lookup.get((src, tgt))
        if edge_id is None:
            continue
        key = f"interf:{edge_id}:{self._INTERF_METRIC}"
        result[key] = self._build_edge_series(
            edge_id, self._INTERF_METRIC, snapshots
        )
    return result
\end{lstlisting}

Il metodo delega tre operazioni al TopologyBuilder: \textit{get\_compliance\_set\_nodes}
restituisce l'insieme dei nodi di $H_{\Phi_i}$; \textit{get\_edges\_for\_compliance\_set}
restituisce le tuple \textit{(source, target)} degli archi interni;
\textit{\_collect\_interference\_edges} itera su tutti gli altri \textit{compliance set} e
raccoglie gli archi di $M^{\text{interf}}$ chiamando \textit{get\_interference\_edges} per ogni
coppia, deduplicando tramite un \textit{set} di tuple già viste. La traduzione da tupla
\textit{(source, target)} a \textit{edge\_id} avviene tramite \textit{\_edge\_id\_lookup}, un
dizionario precalcolato nel costruttore dal file \textit{topology.yaml}. I metodi privati
\textit{\_build\_node\_series} e \textit{\_build\_edge\_series} inseriscono
\textit{float("nan")} per gli snapshot in cui il nodo o l'arco è assente, emettendo un avviso:
il NaN passa ai moduli di \textit{forecasting} che lo gestiscono con \textit{dropna()} prima del
\textit{fitting}.

\subsection{Costante di Dominio per l'Interferenza}
\label{subsec:interf_metric}

La costante di classe \textit{\_INTERF\_METRIC = "throughput\_rps"} è l'unica costante di
dominio non configurabile tramite YAML nell'intero \textit{framework}. La scelta è giustificata
dalla semantica: il \textit{throughput} è l'unico indicatore che quantifica il carico esercitato
da un \textit{compliance set} esterno su una risorsa condivisa. La latenza, al contrario, è una
metrica di effetto, in quanto misura la degradazione ma non quantifica il carico imposto.
Includerla in $M^{\text{interf}}$ introdurrebbe rumore e dipendenze causali spurie.

Il costruttore verifica che questa metrica sia presente in \textit{edge\_metrics} di
\textit{topology.yaml} e, se assente, emette \textit{logger.warning} e imposta il \textit{flag}
\textit{\_interf\_available = False}, rendendo $M^{\text{interf}} = \emptyset$ per tutti i
\textit{compliance set} senza sollevare eccezioni. Questa scelta consente di usare il
\textit{framework} anche con dataset che non dispongono di metriche di \textit{throughput}, a
costo di perdere la capacità di rilevare le interferenze \textit{cross-property}.

% =========================================================
\section{Fase I: Forecasting Locale con Routing Dinamico}
\label{sec:impl_fase1}
% =========================================================

La Fase~I della \textit{pipeline} trasforma le serie temporali di $M_{\Phi_i}$ in previsioni
quantitative sull'orizzonte $\tau$ configurato. L'implementazione di questa fase nel modulo
StatForecaster affronta tre problemi ingegneristici distinti: come selezionare il modello più
adeguato per ogni serie senza intervento manuale, come gestire i \textit{timestamp} irregolari
del dataset, e come garantire che l'addestramento non sia inquinato da dati anomali. Il nucleo
della soluzione ai tre problemi è il metodo \textit{fit()}, che incapsula la cascata di decisioni
e le trasformazioni necessarie.

\subsection{Routing Dinamico dei Modelli}
\label{subsec:routing_strategy}

Lo StatForecaster implementa il \textit{Pattern Strategy} per l'assegnazione dei modelli
predittivi: il comportamento di \textit{fitting} e predizione varia a seconda del modello
selezionato per ogni serie, ma l'interfaccia pubblica (\textit{fit} e \textit{predict}) rimane
invariata indipendentemente dai modelli scelti. La selezione avviene tramite il metodo privato
\textit{\_route\_model}, mostrato nel Codice~\ref{lst:route_model}, che implementa una cascata
di tre regole valutate nell'ordine:

\begin{lstlisting}[
  language=Python,
  caption={Metodo \_route\_model in src/phase1/stat\_forecaster.py.},
  label={lst:route_model}
]
def _route_model(self, key, metric_name, n_samples, model_override):
    if model_override and key in model_override:
        return model_override[key]

    if metric_name in self._nonlinear_metrics:
        if n_samples < self._input_window * 2:
            logger.warning(
                "Serie '%s' ha %d campioni < input_window*2=%d. "
                "Fallback da LSTM a Prophet.",
                key, n_samples, self._input_window * 2,
            )
            return "prophet"
        return "lstm"

    return "prophet"
\end{lstlisting}

La prima regola controlla se la chiave della serie è presente nel dizionario
\textit{model\_override} passato opzionalmente a \textit{fit()}: in caso affermativo, il modello
specificato viene usato incondizionatamente, bypassando ogni altra logica. Questo meccanismo è il
punto di accesso per ARIMA e la regressione lineare, che non entrano mai nel \textit{routing}
automatico: richiedono un \textit{override} esplicito da parte dell'operatore.

La seconda regola verifica se il nome finale della metrica (il segmento dopo l'ultimo carattere
\textit{:} nella chiave) appartiene alla lista \textit{nonlinear\_metrics} letta da
\textit{pipeline\_params["forecasting"]["lstm"]["nonlinear\_metrics"]}. Se sì, il modello di
destinazione è LSTM; ma se il numero di campioni nominali disponibili è inferiore a due volte la
finestra di \textit{input} (\textit{input\_window}), si applica un \textit{fallback} automatico
a Prophet con emissione di avviso. Questo \textit{fallback} è fondamentale per garantire che il
sistema non tenti di addestrare un modello su serie troppo corte tramite feature laggiate, producendo
un \textit{Ridge} con finestre di dimensione zero o negativa.

La terza regola è il \textit{default} incondizionato: Prophet. Questa scelta è motivata dalla
robustezza di Prophet rispetto alle serie temporali di sistemi distribuiti: gestisce nativamente
le irregolarità nei \textit{timestamp}, apprende automaticamente i \textit{changepoint}, non
richiede pre-processing per la stazionarietà e produce intervalli di credibilità bayesiani
calibrati. ARIMA non è nel \textit{routing} automatico perché richiede che l'operatore verifichi
esplicitamente la stazionarietà della serie e selezioni consapevolmente gli ordini; la regressione
lineare non lo è perché la sua utilità è limitata al ruolo di \textit{baseline} di validazione.

Indipendentemente dal percorso seguito dalla cascata, il nome del modello prodotto da
\textit{\_route\_model} viene validato rispetto all'insieme chiuso di modelli ammessi prima
di procedere al \textit{fitting}. Un nome non riconosciuto, come potrebbe essere specificato
per errore in un dizionario \textit{model\_override}, solleva \textit{ValueError} con
messaggio descrittivo che include il nome invalido e l'elenco dei valori ammessi. Senza questa
validazione, un \textit{override} mal formato cadrebbe silenziosamente nel ramo \textit{else}
della catena di \textit{if/elif} e verrebbe addestrato come regressione lineare, senza alcun
segnale dell'errore di configurazione.

\subsection{Metodo fit(): Filtraggio Nominale e Gestione dei Timestamp Irregolari}
\label{subsec:fit_method}

Il metodo \textit{fit()}, mostrato nel Codice~\ref{lst:fit_forecaster}, addestra i modelli
esclusivamente sulle finestre nominali (\textit{label == 0}), identificate dall'insieme di
\textit{timestamp} estratto da \textit{nominal\_snapshots}. Includere le finestre anomale
nell'addestramento produrrebbe modelli che apprendono la distribuzione delle anomalie come
comportamento atteso, invalidando le previsioni future.

\begin{lstlisting}[
  language=Python,
  caption={Metodo fit() in src/phase1/stat\_forecaster.py.},
  label={lst:fit_forecaster}
]
nominal_ts: set[int] | None = None
if nominal_snapshots is not None:
    nominal_ts = {int(s["timestamp"]) for s in nominal_snapshots}

for key, df in features.items():
    metric_name = key.split(":")[-1]

    if nominal_ts is not None:
        train_df = df[df.index.isin(nominal_ts)].dropna()
    else:
        train_df = df.dropna()

    routing = self._route_model(
        key, metric_name, len(train_df), model_override
    )
    self._routing[key] = routing
    self._train_data[key] = train_df

    if routing == "prophet":
        self._models[key] = self._fit_prophet(train_df)
    elif routing == "lstm":
        self._models[key] = self._fit_lstm_placeholder(train_df)
    elif routing == "arima":
        self._models[key] = self._fit_arima(key, train_df)
    else:
        self._models[key] = self._fit_linear(key, train_df)

self._is_fitted = True
\end{lstlisting}

Il \textit{dropna()} applicato al \textit{training set} rimuove i NaN senza interpolazione.
Questa scelta è deliberata: un'interpolazione introdurrebbe valori sintetici che non si trovano
nel dataset originale e che potrebbero influenzare i parametri del modello in modo non
tracciabile. Immediatamente dopo il filtraggio, un controllo esplicito verifica che il
\textit{training set} risultante non sia vuoto: se tutti i \textit{timestamp} nominali non hanno
corrispondenza nella serie temporale della \textit{feature} considerata, il modulo emette un
\textit{logger.warning} e salta la \textit{feature} con un'istruzione \textit{continue}. La
\textit{feature} non sarà presente nel dizionario restituito da \textit{predict()}, e i moduli
superiori che la richiedessero riceverebbero \textit{None} al posto del DataFrame di previsione.
Questo guard è verificato dal test
\textit{test\_fit\_empty\_train\_after\_nominal\_filter\_skips\_gracefully}: il test fornisce
\textit{nominal\_snapshots} con \textit{timestamp} completamente diversi da quelli della
\textit{feature} e verifica che \textit{predict()} non contenga la chiave corrispondente.

La stima della frequenza di campionamento \textit{freq\_us}, necessaria per costruire i
\textit{timestamp} futuri nella predizione, usa la mediana dei diff sull'indice del training set.
La scelta della mediana rispetto alla media è motivata dalla possibile presenza di gap
inter-esperimento molto estesi, che sarebbero \textit{outlier} fortemente destabilizzanti per
una media aritmetica, mentre la mediana li ignora. Il valore \textit{freq\_us} stimato viene
usato per costruire i \textit{timestamp} futuri come
$t_{\text{last}} + (i+1) \times \textit{freq\_us}$, garantendo esattamente
\textit{horizon\_steps} righe coerenti con la finestra di raccolta nominale.

\subsection{I Quattro Modelli Predittivi}
\label{subsec:modelli_predittivi_impl}

Prophet è il modello \textit{default} del \textit{routing} automatico. Il suo addestramento
richiede la conversione dei \textit{timestamp} in µs in \textit{datetime64} tramite
\textit{pd.to\_datetime(index, unit="us")}: questa conversione avviene solo nel metodo privato
\textit{\_fit\_prophet}, non nell'interfaccia pubblica, evitando una trasformazione globale non
necessaria agli altri modelli. Prophet viene addestrato con \textit{changepoint\_prior\_scale=0.05}
per limitare la flessibilità del \textit{trend} e ridurre l'\textit{overfitting} su serie con
pochi dati nominali. L'intervallo di credibilità bayesiano è impostato a
\textit{interval\_width=0.80}, producendo il 10° e il 90° percentile come \textit{yhat\_lower}
e \textit{yhat\_upper}.

ARIMA è attivato solo tramite \textit{model\_override} e seleziona automaticamente l'ordine
$(p, d, q)$ minimizzando il criterio AIC (\textit{Akaike Information Criterion}) con ricerca
esaustiva sullo spazio di ordini definito dai parametri \textit{arima.max\_p},
\textit{arima.max\_d}, \textit{arima.max\_q} letti da \textit{pipeline\_params.yaml}. Le
eccezioni durante il \textit{fitting} di una combinazione $(p, d, q)$ specifica vengono catturate
e ignorate silenziosamente; se nessuna combinazione produce un modello valido, viene sollevato
\textit{RuntimeError} con il nome della \textit{feature}. L'interfaccia SARIMAX di statsmodels è
usata con \textit{enforce\_stationarity=False} e \textit{enforce\_invertibility=False} per
consentire al selettore AIC di esplorare l'intero spazio di ordini senza vincoli preventivi. I
residui del modello ottimale vengono sottoposti al test di Ljung-Box
(\textit{acorr\_ljungbox}) su un numero di ritardi pari a
$\min(10,\lfloor n/5 \rfloor)$: se la \textit{p-value} è inferiore a $0{,}05$, viene
emesso un \textit{logger.warning} con l'ordine $(p^*,d^*,q^*)$ ottimale e la
\textit{p-value} minima. Il test è puramente diagnostico, non altera il modello
selezionato, e segnala i casi in cui la struttura di autocorrelazione residua
non è stata catturata dal modello AIC-ottimale, suggerendo all'operatore di
considerare un ordine superiore o un modello alternativo.

Il modello LSTM è attualmente un \textit{placeholder} computazionale documentato con un
\textit{TODO} nel codice. La classe interna \textit{\_fit\_lstm\_placeholder} addestra un Ridge
Regressor di scikit-learn su \textit{feature} laggiate di larghezza
\textit{min(input\_window, n-1)}, producendo previsioni quasi lineari con incertezza stimata come
margine simmetrico pari al 5\% del valore previsto. Questa soluzione è deliberata: il
\textit{placeholder} consente di verificare l'intera \textit{pipeline} di \textit{routing} e di
integrazione con i moduli superiori senza dipendere da un ciclo di addestramento profondo su GPU.
La sostituzione con un'architettura LSTM PyTorch con stima dell'incertezza tramite \textit{Monte
Carlo Dropout}~\cite{gal2016} è il primo punto di estensione previsto.

La regressione lineare, attivata tramite \textit{model\_override}, stima i coefficienti $a$ e $b$
del modello $\hat{y} = a + b \cdot t$ usando l'indice temporale come regressore. L'incertezza è
calcolata come la deviazione standard dei residui sul \textit{training set}, proiettata come
intervallo $\pm 1{,}96\sigma$ sul futuro.

Il metodo \textit{predict()} costruisce i \textit{timestamp} futuri e delega la generazione
delle previsioni al metodo privato corrispondente al modello addestrato. L'output è un dizionario
che mappa ogni chiave di \textit{feature} a un DataFrame con colonne \textit{yhat},
\textit{yhat\_lower} e \textit{yhat\_upper} e indice \textit{timestamp}. La proprietà
$\textit{yhat\_lower} \leq \textit{yhat} \leq \textit{yhat\_upper}$ è garantita per tutti i
modelli e verificata da \textit{test\_yhat\_lower\_leq\_yhat\_leq\_upper} con tolleranza
$10^{-9}$.

% =========================================================
\section{Fase II: Analisi Causale Guidata dalla Topologia}
\label{sec:impl_fase2}
% =========================================================

Il CausalAnalyzer implementa la Fase~II della \textit{pipeline}, costruendo il grafo causale
orientato che formalizza le dipendenze statisticamente verificate tra le metriche di $M_{\Phi_i}$.
Il modulo applica lo \textit{screening} topologico prima di qualsiasi calcolo statistico:
anziché testare tutte le coppie possibili di metriche (un numero quadratico nel totale delle
\textit{feature}), il CausalAnalyzer categorizza le coppie in base alla loro rilevanza
semantica rispetto al \textit{compliance set} e applica i test statistici solo alle coppie
che hanno motivazione topologica per essere causalmente correlate.

\subsection{Le Quattro Categorie di Coppie Candidate}
\label{subsec:causal_pairs}

Il metodo \textit{\_build\_candidate\_pairs} produce triple \textit{(source\_key, target\_key,
category)} organizzate in tre categorie. Le coppie di categoria \textit{node\_arc} collegano una
\textit{feature} di nodo a una \textit{feature} di arco per cui il nodo è sorgente o
destinazione, a condizione che il nodo non appartenga all'insieme
$\text{Shared}(H_{\Phi_i}, H_{\Phi_j})$: la metrica di stato interno di un componente come causa e la latenza di un arco
uscente dallo stesso componente come effetto. Queste coppie sono le candidate naturali per
relazioni causali dirette tra lo stato interno di un componente e la qualità dell'interazione che
esso serve, e costituiscono il sottoinsieme più rilevante ai fini della \textit{root cause
analysis}. Le coppie di categoria \textit{intra} comprendono tutte le altre coppie tra
\textit{feature} di $M^{\text{direct}}$ non già classificate come \textit{node\_arc}. Le coppie
di categoria \textit{inter} collegano le \textit{feature} di interferenza (\textit{interf:}) alle
metriche dei nodi condivisi: sono le candidate per la verifica della catena causale
\textit{cross-property}.

La deduplica si basa su \textit{frozenset} della coppia di chiavi: nessuna coppia $(A, B)$ e
$(B, A)$ vengono entrambe incluse. Questa scelta ha una conseguenza rilevante: per le coppie
\textit{intra}, la direzione di causalità testata è determinata dall'ordine di enumerazione.
Le chiavi \textit{node:} vengono enumerate prima delle chiavi \textit{edge:}, quindi per una
coppia che coinvolge una metrica di nodo e una di arco viene sistematicamente testata la
direzione nodo come causa e arco come effetto, che è semanticamente corretta. La direzione
inversa non viene mai testata per coppie \textit{intra}. Per le coppie \textit{node\_arc} questa asimmetria
è intenzionale e corretta; per le coppie \textit{intra} generiche è una semplificazione che
potrebbe produrre falsi negativi in scenari dove la propagazione è inversa.

La regola più importante della categorizzazione riguarda il filtro di Pearson: le coppie
\textit{inter} lo bypassano completamente. Questo \textit{bypass}
è giustificato dal fatto che le interferenze \textit{cross-property} si manifestano tipicamente
come effetti soglia non-lineari: la correlazione di Pearson è prossima a zero finché il
\textit{throughput} del flusso interferente rimane sotto la soglia di saturazione, e diventa
significativa solo dopo. Applicare il filtro di Pearson alle coppie \textit{inter} eliminerebbe
sistematicamente le interferenze più pericolose, ossia quelle che si manifestano bruscamente.

Una quarta categoria, \textit{inter2}, identifica le coppie
$(\textit{node:}v\textit{:metric},\ \textit{edge:}e_{\text{int}}\textit{:metric})$
con $v \in \text{Shared}(H_{\Phi_i}, H_{\Phi_j})$: rappresentano il secondo
collegamento della catena \textit{cross-property}, dal nodo condiviso $v$
verso l'arco interno $e_{\text{int}}$.
Come le coppie \textit{inter}, bypassano il filtro di Pearson.

\subsection{La Pipeline a Tre Stadi per Coppia}
\label{subsec:causal_pipeline}

Il primo stadio (\textit{\_align\_series}) interseca gli indici temporali dei due DataFrame,
applica \textit{sort\_values()} sull'insieme comune per garantire l'ordine temporale crescente,
e poi \textit{dropna()} sull'intersezione ordinata. Il riordinamento è necessario perché le
operazioni di join e filtraggio sui DataFrame pandas non garantiscono la conservazione dell'ordine
temporale dell'indice; i test di Granger e di Transfer Entropy assumono implicitamente che le
osservazioni siano ordinate nel tempo, e serie disordinate producono stime dei \textit{lag}
semanticamente errate. Se il numero di campioni risultante è inferiore a tre, la coppia viene
scartata con \textit{logger.warning}: tre campioni sono il minimo assoluto per calcolare una
statistica significativa.

Il secondo stadio è il test di Granger, preceduto dal pre-processing \gls{ADF}. Il metodo
\textit{\_make\_stationary\_pair} applica il test \gls{ADF} su entrambe le serie effetto e
causa indipendentemente; il numero di differenziazioni applicato è il massimo tra quelli
necessari alle due serie (al massimo due), garantendo la stazionarietà di entrambe prima del
test di Granger. Differenziare entrambe le serie in modo sincrono è essenziale per preservare
la relazione temporale tra causa ed effetto. Se dopo 2 differenziazioni la stazionarietà non
è raggiunta, viene emesso un avviso e l'analisi prosegue. Il numero massimo di
differenziazioni (2) e la soglia ADF (0,05) non sono configurabili da
\textit{pipeline\_params.yaml}: corrispondono per coerenza ai parametri \textit{max\_d} e alla
soglia standard adottati nel modello ARIMA della Fase~I. Il Codice~\ref{lst:granger_te} mostra
il metodo complessivo di analisi per coppia.

\begin{lstlisting}[
  language=Python,
  caption={Metodo \_run\_granger\_then\_te in src/phase2/causal\_analyzer.py.},
  label={lst:granger_te}
]
def _run_granger_then_te(self, source_key, target_key, cause, effect):
    granger_res = self._granger_test(
        cause, effect,
        self._granger_max_lag, self._granger_significance
    )
    if granger_res is not None:
        return {
            "source": source_key, "target": target_key,
            "type": "linear",
            "intensity": granger_res["intensity"],
            "method": "granger",
            "lag": granger_res["lag"],
        }

    te_val = self._transfer_entropy(cause, effect, self._n_bins)
    if te_val > self._te_threshold:
        return {
            "source": source_key, "target": target_key,
            "type": "nonlinear",
            "intensity": float(te_val),
            "method": "transfer_entropy",
            "lag": None,
        }

    return None
\end{lstlisting}

La funzione \textit{\_granger\_test} usa \textit{grangercausalitytests} di statsmodels con
l'output a console soppresso tramite \textit{contextlib.redirect\_stdout}. Il \textit{lag}
ottimale è quello con il $p$-value minimo sul test $F$ (\textit{ssr\_ftest}). L'intensità è
calcolata come $\Delta R^2 = \max(0{,}0,\, R^2_{\text{completo}} - R^2_{\text{ristretto}})$:
il \textit{clipping} a zero previene valori negativi da arrotondamento numerico nella
sottrazione in virgola mobile. La funzione restituisce \textit{None} se il $p$-value migliore
supera la soglia di significatività o se i dati sono insufficienti ($n < \textit{max\_lag} + 2$).

Il terzo stadio, la Transfer Entropy, è riservato alle coppie che non superano Granger. 
La stima usa un approccio discreto: entrambe le serie vengono discretizzate uniformemente in
\textit{n\_bins} classi e il valore della TE è normalizzato sull'entropia di Shannon $H(Y)$
della serie effetto, producendo un valore in $[0, 1]$ confrontabile tra coppie diverse.
Il parametro \textit{n\_bins} (valore di default 10) è dichiarato esplicitamente in
\textit{pipeline\_params.yaml} nella sezione \textit{causal\_analysis} ed è quindi
configurabile senza modificare il codice. Valori più alti riducono il bias campionario
della stima discreta ma richiedono un numero di campioni crescente con $n_{\text{bins}}^2$
per produrre stime affidabili. Il valore 0.0 è restituito quando $H(Y) \leq 0$ (serie costante, entropia nulla) 
o quando i dati sono insufficienti.

\subsection{Le Catene Causali Cross-Property}
\label{subsec:cross_property_impl}

Il metodo \textit{\_check\_cross\_property} cerca strutture causali della forma
$\textit{interf:e\_j} \to \textit{node:v} \to \textit{edge:e\_internal}$ dove
$v \in \text{Shared}(H_i, H_j)$ e $\textit{e\_internal} \in A(H_{\Phi_i})$. Per ogni catena
candidata vengono eseguiti due test di Granger consecutivi: il campo \textit{confirmed} è
\textit{True} solo se entrambi superano la soglia di significatività. Le coppie auto-referenziali,
dove l'\textit{edge\_id} della \textit{feature} di interferenza coincide con quello dell'arco
interno, vengono saltate per evitare test triviali. Le eccezioni per singola coppia o singola
catena sono catturate e registrate come avvisi senza interrompere la \textit{pipeline}: il modulo
restituisce sempre un grafo causale valido, anche parziale.

% =========================================================
\section{Fase III: Monitoraggio Strutturale Gerarchico}
\label{sec:impl_fase3}
% =========================================================

Lo StructuralMonitor è l'unico modulo con stato (\textit{stateful}) dell'intero \textit{framework}.
Mantiene l'accumulatore \gls{CUSUM} e lo stato \gls{EWMA} tra chiamate successive a
\textit{monitor()}, consentendo il rilevamento di derive persistenti che attraversano molteplici
finestre temporali consecutive. Tutti gli altri moduli della \textit{pipeline} sono privi di
stato (\textit{stateless}) e producono lo stesso output a parità di \textit{input}.
L'implicazione pratica è che \textit{reset\_cusum()} deve essere chiamato esplicitamente tra
esperimenti distinti: il \textit{carry-over} dello stato di un esperimento nel successivo
produrrebbe accumulatori già carichi all'inizio del nuovo esperimento, generando falsi positivi
immediati.

\subsection{Metodo fit(): Calibrazione sui Dati Nominali}
\label{subsec:monitor_fit}

Il metodo \textit{fit()} addestra il monitor in cinque passi sequenziali che devono essere
eseguiti nell'ordine specificato, perché ciascuno dipende dai risultati del precedente. Il primo
passo costruisce il vettore di stato aggregato $\mathbf{X}_{H_{\Phi_i}}(t)$ concatenando le
\textit{feature} di nodo di tutti i componenti del \textit{compliance set}, nell'ordine
deterministico \textit{sorted(cs\_nodes) $\times$ sorted(node\_metrics)}. I valori NaN nel
vettore aggregato sono imputati con la media di colonna calcolata sui soli snapshot nominali.
Le \textit{feature} di arco sono escluse perché effimere: nelle finestre con traffico assente
i loro valori sarebbero strutturalmente nulli indipendentemente dallo stato di salute del sistema,
rendendo instabile l'addestramento dell'Isolation Forest.

Il secondo passo addestra l'Isolation Forest sul vettore aggregato con i parametri letti da
\textit{pipeline\_params["anomaly\_detection"]["isolation\_forest"]}. Il terzo passo calcola
la media $\mu_m$ e la deviazione standard $\sigma_m$ nominali per ogni \textit{feature} di
$M_{\Phi_i}$, necessarie per il calcolo adattivo dello z-score. Il quarto passo determina il
valore di riferimento per il \gls{CUSUM}: per le topologie lineari calcola
$\text{PAS}_{\text{gold}}$ invocando \textit{compute\_path\_adherence} del PBOBuilder con una
\textit{weight series} sintetica di un solo elemento costruita da $W^{\text{gold}}$; per le
topologie parallele imposta il riferimento a zero, perché $\text{Frobenius}_{\text{gold}} = 0$
per costruzione (il \textit{gold standard} è il comportamento nominale). Il quinto passo chiama
\textit{reset\_cusum()}, azzerando l'accumulatore, lo stato \gls{EWMA} e la cronologia della
\textit{deque}.

\subsection{Metodo monitor(): Gerarchia di Attivazione a Quattro Livelli}
\label{subsec:monitor_levels}

Il metodo \textit{monitor()} esegue i quattro livelli della gerarchia in una sequenza che
rispecchia direttamente la logica descritta nel Capitolo~\ref{chap:metodologia},
come mostrato nel Codice~\ref{lst:monitor_hierarchy}:

\begin{lstlisting}[
  language=Python,
  caption={Gerarchia di attivazione in monitor() di src/phase3/structural\_monitor.py.},
  label={lst:monitor_hierarchy}
]
# Livello 1a -- threshold SLA
threshold_violations = self._check_threshold(
    compliance_set_name, features, timestamp
)
# Livello 1b -- z-score adattivo
zscore_violations = self._check_zscore(features, timestamp)
base_signal = (
    len(threshold_violations) > 0 or len(zscore_violations) > 0
)
# Livello 2 -- Isolation Forest (condizionale a base_signal)
if base_signal:
    if_signal = self._check_isolation_forest(features, timestamp)
else:
    if_signal = False
# Livello 3 -- EWMA + CUSUM (sempre, in parallelo)
(cusum_signal, cusum_stat, ewma_value,
 frobenius_distance, pas_value) = self._update_ewma_cusum(
     weight_series, compliance_set_name
)
# Livello 4 -- Validatore strutturale (solo se IF e CUSUM entrambi attivi)
structural_confirmed = False
if if_signal and cusum_signal:
    structural_confirmed = self._check_structural_validator(
        frobenius_distance, pas_value
    )
\end{lstlisting}

I Livelli~1a e~1b sono sempre attivi e costituiscono il filtro deterministico di base. Il
Livello~1a confronta il valore corrente di ogni \textit{feature} con la soglia \gls{SLA}
corrispondente letta da \textit{topology.yaml}: il campo \textit{bound} determina la direzione
(\textit{upper} per latenza ed error rate, \textit{lower} per \textit{throughput}). I valori NaN
non sono mai considerati violazioni. Il Livello~1b calcola lo z-score adattivo
$z_m(t) = (m(t) - \mu_m) / \sigma_m$ per ogni \textit{feature}: la soglia
\textit{zscore\_threshold} (di default 3.0, corrispondente allo 0,3\% delle osservazioni sotto
una gaussiana) è letta da \textit{pipeline\_params}. Il segnale \textit{base\_signal} è vero se
almeno una \textit{feature} viola una delle due condizioni.

Il Livello~2 applica l'Isolation Forest sul vettore di stato aggregato costruito per lo snapshot
corrente, con la stessa struttura e lo stesso ordine usati durante il \textit{fit()}. I valori
NaN vengono imputati con la media di colonna del \textit{training} prima della classificazione.
Il Livello~2 è attivato solo se \textit{base\_signal} è vero: questa condizione riduce il carico
computazionale evitando la classificazione IF nelle finestre in cui non vi è alcun segnale metrico
di base. Il gating introduce tuttavia una limitazione: anomalie puramente multivariate, in cui
ciascuna metrica rimane individualmente entro le soglie \gls{SLA} e lo z-score, ma la loro
combinazione è anomala rispetto al profilo storico, non vengono mai sottoposte all'Isolation
Forest perché \textit{base\_signal} è falso. Questo scenario è distinto da quello in cui una
singola metrica è fuori soglia e IF conferma l'anomalia aggregata; richiederebbe un'architettura
alternativa in cui IF opera su tutte le finestre indipendentemente dai segnali puntuali, con costi
computazionali maggiori.

Il Livello~3 opera in parallelo rispetto ai Livelli~1 e~2: non dipende dal valore di
\textit{base\_signal} e non ne è condizionato. Il \gls{PAS} o la norma di Frobenius corrente
vengono calcolati tramite il PBOBuilder, filtrati dall'\gls{EWMA} e accumulati nel \gls{CUSUM}.
La scelta della metrica dipende dal \textit{topology\_type}: per le topologie lineari si accumula
il decremento del \gls{PAS}, per le parallele l'incremento della norma di Frobenius. Il
Codice~\ref{lst:ewma_cusum} mostra l'implementazione del filtro e dell'aggiornamento.

\begin{lstlisting}[
  language=Python,
  caption={Filtro EWMA e aggiornamento CUSUM in src/phase3/structural\_monitor.py.},
  label={lst:ewma_cusum}
]
def _apply_ewma(self, signal_raw: float) -> float:
    if self._ewma_state is None:
        self._ewma_state = signal_raw   # cold start
    else:
        self._ewma_state = (
            self._ewma_alpha * signal_raw
            + (1 - self._ewma_alpha) * self._ewma_state
        )
    self._ewma_history.append(self._ewma_state)
    return self._ewma_state

# Aggiornamento CUSUM (topologia lineare):
ewma_new = self._apply_ewma(pas_value)
increment = self._reference - ewma_new - self._cusum_k
self._cusum_stat = max(0.0, self._cusum_stat + increment)
cusum_signal = self._cusum_stat > self._cusum_threshold
\end{lstlisting}

Il \textit{cold start} inizializza lo stato \gls{EWMA} con la prima osservazione disponibile,
evitando distorsioni nelle prime finestre. L'accumulatore \gls{CUSUM} accumula il termine
$\text{PAS}_{\text{gold}} - \tilde{s}_t - k$: cresce in presenza di degrado e decresce
naturalmente durante il recovery. Il \textit{clipping} a zero garantisce la non-negatività
di \textit{cusum\_stat}. Il metodo \textit{reset\_cusum()} va invocato esplicitamente tra
sessioni distinte per azzerare i residui accumulati.

Il Livello~4, il Validatore strutturale, è attivato solo quando sia \textit{if\_signal} sia
\textit{cusum\_signal} sono veri, e verifica due condizioni congiunte. La prima riguarda la distanza strutturale
attuale: per le topologie lineari questa corrisponde al
PAS-gap $|\text{PA}(t) - \text{PAS}_{\text{gold}}| > \delta$
come descritto in~\cite[§3.2.3]{methodology}, e deve superare \textit{distance\_threshold} (default 0,15);
per le topologie parallele corrisponde alla norma di Frobenius
$\|W(t) - W^{\text{gold}}\|_F > \delta$ con la stessa soglia.
La seconda condizione richiede che gli ultimi \textit{consecutive\_windows} incrementi consecutivi
della \textit{deque} \gls{EWMA} abbiano il segno di degrado corretto (negativo per il \gls{PAS},
positivo per la Frobenius). La persistenza della deriva per almeno
\textit{consecutive\_windows} finestre consecutive è il discriminante tra una fluttuazione
momentanea e un allontanamento strutturale autentico.

% =========================================================
\section{Fase IV: Generazione dell'Alert}
\label{sec:impl_fase4}
% =========================================================

L'AlertGenerator implementa la Fase~IV come un'operazione di sintesi pura: non esegue calcoli
statistici propri, ma aggrega i risultati delle tre fasi precedenti (previsioni, grafo causale,
risultato del monitoraggio) per produrre un giudizio di \textit{compliance} predittivo
strutturato. Il metodo principale \textit{generate()} produce un dizionario con quindici campi
che costituisce l'\textit{alert} strutturato, oppure restituisce \textit{None} se nessuna
violazione è prevista nell'orizzonte considerato. La restituzione di \textit{None} non è un
errore: è la rappresentazione esplicita della conformità predittiva del sistema nell'orizzonte
corrente.

\subsection{Proprietà Monitorata e Funzione di Aggregazione}
\label{subsec:property_detection}

La prima operazione del metodo \textit{generate()} è identificare quale proprietà sia monitorata
per il \textit{compliance set} specificato, in modo da selezionare la funzione di aggregazione
appropriata. Il metodo \textit{\_detect\_property} ispeziona la sezione \textit{sla} del
\textit{compliance set} in \textit{topology.yaml} con priorità decrescente: latenza, affidabilità,
capacità. Per ognuna, determina la soglia di controllo interna e la soglia di visualizzazione
nell'\textit{alert}, che possono differire. Per la proprietà di affidabilità, la \gls{SLA} è
espressa come \textit{error\_rate upper 0.05}, ma il controllo interno opera sulla
\textit{reliability} $\hat{R} = \prod_e (1 - \hat{\varepsilon}_e)$: il modulo converte
automaticamente la soglia da \textit{error\_rate upper 0.05} a \textit{reliability lower 0.95}
per il \textit{violation check}, riportando nell'\textit{alert} i valori originali del certificato
per leggibilità dell'operatore.

La funzione di aggregazione è determinata congiuntamente dalla proprietà e dal
\textit{topology\_type}: per la latenza su topologia lineare si usa la somma, su topologia
parallela il massimo. Per l'affidabilità si usa il prodotto dei complementi $\hat{R}$. Per la capacità il minimo.

Quando il valore previsto per una \textit{feature} è NaN, il modulo applica un
\textit{fallback} conservativo: il valore sostitutivo è scelto in modo da non mascherare
una violazione. La scelta si basa sulla direzione della soglia così come è espressa nel
certificato originale (\textit{display bound}), non su quella risultante dall'eventuale
conversione interna. Per le proprietà con soglia superiore il \textit{fallback} è la
soglia \gls{SLA} stessa, che colloca la previsione al margine del vincolo. Per la
capacità, con soglia inferiore, il \textit{fallback} è 0,0. Questa distinzione è
necessaria per la proprietà di affidabilità: se si usasse la direzione interna
convertita, il \textit{fallback} produrrebbe un tasso di errore nullo e una stima
di affidabilità pari a 1,0, ossia l'esito opposto di quello conservativo cercato.

\subsection{Lead Time e Classificazione della Criticità}
\label{subsec:lead_time_criticality}

Il \textit{lead time} $\tau^*$ è il minimo indice $\tau' \in \{1, \ldots, \tau\}$ per cui
l'aggregato previsto viola la condizione \gls{SLA}. Il metodo \textit{\_estimate\_lead\_time}
itera sull'array delle previsioni aggregate e confronta ogni valore con la soglia: per
\textit{bound} superiore controlla $\hat{\Phi}_i(t+\tau') > T_m$, per \textit{bound} inferiore
$\hat{\Phi}_i(t+\tau') < T_m$. I valori NaN nell'aggregato non generano violazione. Se
l'insieme delle violazioni è vuoto, \textit{generate()} restituisce \textit{None}.

La classificazione della criticità implementa la cascata descritta nel
Capitolo~\ref{chap:metodologia}, come mostrato nel Codice~\ref{lst:classify_criticality}:

\begin{lstlisting}[
  language=Python,
  caption={Metodo \_classify\_criticality in src/phase4/alert\_generator.py.},
  label={lst:classify_criticality}
]
def _classify_criticality(self, lead_time_steps, monitor_result):
    lead_time_days = lead_time_steps * self._step_duration_hours / 24.0

    cusum  = bool(monitor_result.get("cusum_signal", False))
    if_sig = bool(monitor_result.get("if_signal", False))
    confirmed = bool(monitor_result.get("structural_confirmed", False))

    if (
        lead_time_days < self._orange_min_days
        or (cusum and confirmed)
        or (if_sig and confirmed)
    ):
        return "red"

    if (
        self._orange_min_days <= lead_time_days < self._yellow_min_days
        or cusum
        or if_sig
    ):
        return "orange"

    return "yellow"
\end{lstlisting}

Il livello Rosso si attiva quando il \textit{lead time} è inferiore a \textit{orange\_min\_days}
(di default 2 giorni), oppure quando il \gls{CUSUM} e il Validatore strutturale sono entrambi
attivi, oppure quando l'Isolation Forest e il Validatore strutturale sono entrambi attivi. Le
ultime due condizioni segnalano che la deviazione è confermata sia sul piano metrico aggregato
sia su quello comportamentale strutturale. Il livello Arancione copre i casi intermedi: un
\textit{lead time} tra 2 e 7 giorni, oppure uno dei due segnali (\gls{CUSUM} o Isolation Forest)
attivo senza conferma strutturale. Il livello Giallo è il residuale.

Un vincolo implicito di configurazione riguarda la relazione tra le due soglie: il parametro
\textit{orange\_min\_days} deve essere strettamente inferiore a \textit{yellow\_min\_days}.
Se i valori fossero invertiti, la cascata di regole produrrebbe classificazioni semanticamente
errate senza sollevare eccezioni, poiché il costruttore non valida questa relazione. La
correttezza della configurazione è demandata al file \textit{pipeline\_params.yaml}.
Dopo la classificazione primaria, il metodo applica il meccanismo di \textit{declassamento per
incertezza modellistica}: se il \textit{flag} \textit{model\_uncertainty\_flag} è attivo (la
divergenza tra previsione principale e \textit{baseline} lineare supera
\textit{divergence\_threshold}) e il livello corrente non è già Giallo, il livello viene abbassato
di un gradino (da Rosso ad Arancione, da Arancione a Giallo). Il Giallo non viene ulteriormente
declassato. Questo meccanismo garantisce che un \textit{alert} Rosso basato su una previsione
fortemente non-lineare venga automaticamente declassato ad Arancione, riducendo il rischio di
interventi tecnici urgenti motivati da previsioni a bassa affidabilità.

\subsection{Struttura dell'Alert ed Estrazione della Causa Radice}
\label{subsec:alert_structure}

Il dizionario \textit{alert} prodotto da \textit{generate()} contiene quindici campi. I campi
principali sono \textit{timestamp}, \textit{compliance\_set}, \textit{property\_at\_risk} e
\textit{criticality}. Il \textit{lead\_time\_steps} indica il numero di passi temporali fino alla
prima violazione prevista, mentre \textit{lead\_time\_hours} è il suo equivalente in ore,
ottenuto moltiplicando per \textit{step\_duration\_hours} letto da \textit{pipeline\_params}.
Il campo \textit{aggregated\_forecast} contiene il vettore delle previsioni aggregate su
tutti gli step dell'orizzonte; \textit{sla\_threshold} e \textit{sla\_bound} riportano i valori
originali del certificato, distinti dai valori interni di controllo che possono differire per la
proprietà di affidabilità. Il campo \textit{structural\_signals} è un sotto-dizionario che riporta
tutti i segnali prodotti dallo StructuralMonitor: \textit{base\_signal}, \textit{if\_signal},
\textit{cusum\_signal}, \textit{structural\_confirmed}, \textit{frobenius\_distance} e
\textit{pas\_value}.

L'estrazione della causa radice analizza il grafo causale di Fase~II in tre passi. Il primo
filtra gli archi causali il cui \textit{target} termina con il suffisso della metrica a
rischio (\textit{:latency\_ms}, \textit{:error\_rate} o \textit{:throughput\_rps}).
Il secondo seleziona l'arco con intensità massima nell'insieme filtrato: la \textit{source} di
quell'arco è la causa radice, con due livelli di \textit{fallback} innestati. Se l'insieme
filtrato è vuoto, la ricerca si restringe ai soli archi il cui campo \textit{target} inizia
con \textit{"edge:"}, escludendo le coppie tra metriche di nodo prive di correlazione
diretta con un arco topologico. Se anche questo insieme è vuoto, oppure se il campo
\textit{target} dell'arco selezionato non è di tipo \textit{"edge:"}, si attiva il
\textit{fallback} finale: \textit{\_find\_critical\_arc\_from\_forecast} identifica l'arco
critico come quello con il valore previsto più estremo al passo del \textit{lead time},
senza dipendere dal grafo causale. Il terzo passo cerca tra le catene \textit{cross-property}
la prima con \textit{confirmed == True} e ne riporta il \textit{source\_cs} come
\textit{cross\_property\_interference} e la sequenza di nodi come \textit{causal\_chain}.

% =========================================================
\section{Sviluppo Guidato dai Test}
\label{sec:tdd}
% =========================================================

L'adozione sistematica del \gls{TDD} ha prodotto una suite di test con caratteristiche
architetturali precise, che riflettono direttamente i vincoli di progettazione descritti nelle
sezioni precedenti.

Tutti i test operano esclusivamente con strutture dati costruite in memoria. Nessun test legge un
file dal disco, salvo i test del ConfigLoader che usano file temporanei costruiti su
\textit{tmp\_path}. Questo rende la suite eseguibile in pochi secondi anche su macchine senza
accesso ai dataset di produzione. Il modello a dipendenze unidirezionali tra i pacchetti è
verificato implicitamente dalla struttura dei test: un test del TopologyBuilder che importasse
classi del StatForecaster violerebbe la gerarchia e produrrebbe un'importazione ciclica immediata.

Un principio fondamentale dell'intera suite è la verifica dell'assenza di \textit{hardcoding}: i
test calcolano i valori attesi direttamente dalla configurazione caricata, non da costanti fisse
nel codice. Il Codice~\ref{lst:test_hcrit}, tratto dalla suite del FeatureSelector, ne fornisce
un esempio rappresentativo.

\begin{lstlisting}[
  language=Python,
  caption={Test test\_h\_crit\_direct\_node\_count in tests/test\_feature\_selector.py.},
  label={lst:test_hcrit}
]
def test_h_crit_direct_node_count(
    feature_selector, topology_builder, config, mock_snapshots
):
    topo = config.load_topology()
    expected = (
        len(topology_builder.get_compliance_set_nodes("H_crit"))
        * len(topo["node_metrics"])
    )
    result = feature_selector.select_features("H_crit", mock_snapshots)
    node_keys = [k for k in result if k.startswith("node:")]
    assert len(node_keys) == expected
\end{lstlisting}

Il test non usa una costante numerica fissa: interroga il TopologyBuilder
per il numero di nodi del \textit{compliance set} e il ConfigLoader per il numero di metriche,
moltiplicandoli. Se \textit{topology.yaml} evolvesse aggiungendo una nuova metrica, il valore atteso si
aggiornerebbe automaticamente senza richiedere alcuna modifica al codice di test. La suite
rimane valida per costruzione rispetto alle evoluzioni della configurazione.

La suite copre complessivamente 296 casi distribuiti in dieci file. I tipi di test
includono la struttura degli output (forma, chiavi, tipi attesi), la correttezza numerica con
tolleranze esplicite ($\varepsilon = 10^{-9}$ per i confronti in virgola mobile), la gestione dei
casi limite (NaN, liste vuote, \textit{throughput} nullo, valori negativi, JSON malformato), il
comportamento degli errori verificati con \textit{pytest.raises(ExcType, match="...")} per
controllare anche il contenuto del messaggio, e l'assenza di effetti collaterali come
l'idempotenza di \textit{build()}, l'isolamento della \textit{cache} e la non-corruzione dello
stato tra test consecutivi.

% =========================================================
\section{Architettura Online e Deployment}
\label{sec:architettura_online}
% =========================================================

Il \textit{framework} è progettato per essere operato in due fasi temporalmente distinte: una
fase di addestramento \textit{offline}, eseguita una volta su dati storici nominali, e un ciclo
di \textit{inference} \textit{online}, eseguito continuamente su ogni nuova finestra di
telemetria. La distinzione è fondamentale per comprendere quali componenti mantengono stato
tra un'esecuzione e l'altra e quali sono invece privi di stato e rieseguibili in modo
idempotente.

\subsection{Separazione tra Addestramento e Inferenza}
\label{subsec:training_inference}

Nella fase di addestramento, tutti i moduli che richiedono una calibrazione su dati storici
vengono invocati sui soli snapshot nominali. Lo StatForecaster viene addestrato sulle
serie temporali nominali, apprende i parametri dei modelli predittivi e stima la frequenza
di campionamento. Il CausalAnalyzer produce il grafo causale a partire dalle dipendenze
statistiche estratte dalla storia nominale. Lo StructuralMonitor calibra le statistiche
nominali per lo z-score, addestra l'Isolation Forest e calcola i valori di riferimento per il
\gls{CUSUM}. Il PBOBuilder calcola il \textit{gold standard} $W^{\text{gold}}$ dalla
distribuzione media del traffico nelle finestre nominali. Al termine di questa fase, lo stato
addestrato di ciascun modulo viene serializzato per essere riutilizzato nelle sessioni di
\textit{inference} successive.

Nella fase di \textit{inference}, il ciclo operativo elabora ogni nuova finestra di metriche
prodotta dall'infrastruttura di osservabilità. Per ogni finestra, il FeatureSelector
estrae le serie temporali rilevanti per il \textit{compliance set} monitorato, lo
StatForecaster genera le previsioni sull'orizzonte $\tau$ configurato, lo
StructuralMonitor aggiorna il proprio stato interno (\gls{CUSUM} e \gls{EWMA}) e
classifica lo snapshot corrente, e infine l'AlertGenerator produce il giudizio di
\textit{compliance} predittivo. I risultati del CausalAnalyzer, che vengono ricalcolati
\textit{offline} su base periodica (ad esempio con cadenza giornaliera su una finestra scorrevole
di dati nominali recenti), vengono passati all'AlertGenerator come struttura pre-calcolata.

Tre aspetti operativi meritano attenzione nel contesto di un deployment continuo. Lo
StatForecaster richiede un riaddestramento periodico per adattarsi al \textit{behavioral
drift} del sistema: Prophet può essere riaddestrato con cadenza settimanale su una finestra
scorrevole di dati nominali recenti senza interrompere il ciclo di \textit{inference}. 
Lo StructuralMonitor, in quanto unico modulo stateful della pipeline, richiede che
\textit{reset\_cusum()} venga invocato esplicitamente dall'orchestratore quando
il sistema è tornato nominale dopo un episodio di degrado, per azzerare l'accumulatore. Il
PBOBuilder e il relativo \textit{gold standard} devono essere ricalibrati ogni volta
che il sistema subisce modifiche architetturali significative o variazioni strutturali nei
pattern di traffico.

\subsection{Punti di Estensione verso la Produzione}
\label{subsec:estensioni_produzione}

L'implementazione attuale copre l'intera \textit{pipeline} di elaborazione dalla telemetria
agli \textit{alert}. Il componente mancante per un deployment in produzione è l'orchestratore
che governa il ciclo di acquisizione e invoca la \textit{pipeline} a ogni finestra temporale.
La funzione di conversione del singolo snapshot dovrebbe essere implementata come
\textit{convert\_window()} nel modulo \gls{ETL}, con la stessa interfaccia contrattuale di
\textit{convert\_file()} ma ottimizzata per la latenza di elaborazione incrementale anziché
per il \textit{throughput} del batch.

Per estendere il sistema a una sorgente diversa, come Prometheus o OpenTelemetry in tempo reale,
è sufficiente implementare un PrometheusConverter con la medesima interfaccia del DSBConverter:
produrre i tre DataFrame canonici (\textit{node\_metrics}, \textit{edge\_metrics},
\textit{ground\_truth}) a partire dalle metriche raccolte nell'ultima finestra temporale.
Tutto il codice a valle della conversione, dai moduli ATGBuilder fino all'AlertGenerator,
funzionerebbe senza modifiche, per costruzione dell'architettura a formato canonico.